\documentclass[11pt]{article}
\usepackage[margin=1in]{geometry}
\usepackage{amsmath,amssymb,amsthm}
\usepackage[hidelinks]{hyperref}
\newtheorem{theorem}{Theorem}
\newtheorem{lemma}{Lemma}
\title{An exact counterexample to the generalized Lambert inverse-radius formula}
\author{Principia Math In Progress}
\date{18 September 2026}
\begin{document}
\maketitle
\begin{abstract}
For \(f(z)=z\sqrt{1+z}\,e^z\), the inverse germ at zero has Taylor
radius \(1/(2\sqrt{2e})\). The two values permitted by Castle's
Conjecture 1.2 are instead \(1/(2\sqrt2\,e)\) and \(2/e\).
All integer phase choices have modulus one, so none repairs the
discrepancy. We also prove an exact radius theorem for arbitrary
positive real powers of factors with negative real roots.
\end{abstract}

\section*{Task 1: the original conjecture}
Let
\[
 f(z)=z\prod_{j=1}^m(z-t_j)^{p_j}e^z,\qquad
 t_j,p_j\in\mathbb C\setminus\{0\}.
\]
Use the primary source's specified branches near zero, where the
inverse germ exists whenever these branches are analytic.
Conjecture 1.2 of \cite[p.~3]{Castle}, reproduced by MathDB 338324,
asserts that the Taylor radius of this inverse equals
\begin{equation}\label{claim}
 \left|
 e^{\,2\pi i\sum_j l_j\phi_j-1}
 \left(1-\sum_j\phi_j\right)^{1-\sum_j\phi_j}
 \prod_j
 \left(\frac{t_j}{p_j+\phi_j}\right)^{p_j+\phi_j}
 p_j^{p_j}\phi_j^{\phi_j}
 \right|
\end{equation}
for some integers \(l_j\) and some solution \(\phi\) of
\begin{equation}\label{saddles}
 \left(1-\sum_j\phi_j\right)(p_i+\phi_i)+t_i\phi_i=0
 \quad(1\le i\le m).
\end{equation}
The exponential in \eqref{claim} is copied from the primary
preprint, including its real term \(-1\).

\begin{theorem}\label{counter}
Take \(m=1,\ t_1=-1,\ p_1=1/2\), and the square root that equals
one at \(z=0\). For
\[
 f(z)=z\sqrt{1+z}\,e^z
\]
the inverse germ \(g(0)=0\) has Taylor radius
\[
 R=\frac{e^{-1/2}}{2\sqrt2}.
\]
All choices in \eqref{claim} give one of the two different values
\[
 R_+=\frac{e^{-1}}{2\sqrt2},\qquad R_-=\frac2e.
\]
Thus Conjecture 1.2 is false.
\end{theorem}

\subsection*{Exact location of the radius}
We first prove a general elementary statement that excludes
unseen complex singularities closer to the origin.

\begin{lemma}\label{positive}
Let \(m\ge1\), \(a_j>0\), \(p_j>0\), and
\[
 f(z)=ze^z\prod_{j=1}^m(a_j+z)^{p_j},
\]
with the powers positive at zero. There is a unique
\(\tau\in(0,\min_j a_j)\) satisfying
\begin{equation}\label{tau}
 1-\tau-\sum_{j=1}^m\frac{p_j\tau}{a_j-\tau}=0.
\end{equation}
It also satisfies \(\tau<1\). The inverse germ of \(f\) at zero
has Taylor radius
\begin{equation}\label{radius}
 \rho=\tau e^{-\tau}\prod_{j=1}^m(a_j-\tau)^{p_j}.
\end{equation}
\end{lemma}

\begin{proof}
Put \(a=\min_j a_j\) and, for \(|u|<a\), define
\[
 \Phi(u)=e^u\prod_j(a_j-u)^{-p_j},\qquad
 F(u)=\frac{u}{\Phi(u)}=-f(-u).
\]
Every Taylor coefficient of \(\Phi\) is positive. This follows by
multiplying \(e^u\) and the binomial expansions
\[
 (a_j-u)^{-p_j}
   =a_j^{-p_j}\sum_{k\ge0}\frac{(p_j)_k}{k!}
                       \left(\frac{u}{a_j}\right)^k.
\]
The function
\[
 L(u)=1-u-\sum_j\frac{p_ju}{a_j-u}
\]
strictly decreases on \((0,a)\), from \(1\) to \(-\infty\), because
\[
 L'(u)=-1-\sum_j\frac{p_ja_j}{(a_j-u)^2}<0.
\]
This gives the unique root \(\tau\). If \(a>1\), then \(L(1)<0\);
otherwise \(\tau<a\le1\). Hence always \(\tau<1\).
Since \(F'(u)=F(u)L(u)/u\), the function \(F\) increases strictly
from \(0\) to \(F(\tau)=\rho\) on \([0,\tau]\), with \(F'(\tau)=0\).

Let \(h(w)=-g(-w)\) be its inverse germ. Lagrange inversion gives
\begin{equation}\label{lagrange}
 h(w)=\sum_{n\ge1}c_nw^n,\qquad
 c_n=\frac1n[u^{n-1}]\Phi(u)^n>0.
\end{equation}
This formula also follows directly by changing variables \(w=F(u)\)
in the coefficient residue integral and integrating by parts.
For any \(0<r<a\), positivity of coefficients gives
\[
 c_n r^{n-1}\le\frac{\Phi(r)^n}{n},
 \qquad c_n\le\frac{r}{nF(r)^n}.
\]
In particular, taking \(r=\tau\) proves convergence for every
\(|w|<\rho\). Thus the radius is at least \(\rho\).

For real \(0\le w<\rho\), the series agrees with the increasing
real inverse of \(F\) on \([0,\tau)\): both agree near zero and
continue uniquely along this interval since \(F'>0\) there.
Consequently \(h(w)\to\tau\) as \(w\uparrow\rho\).
If the Taylor radius exceeded \(\rho\), then \(h\) would be
analytic at \(\rho\), with \(h(\rho)=\tau\). The identity
\(F(h(w))=w\) would then imply
\[
 F'(\tau)h'(\rho)=1,
\]
which is impossible. The radius is therefore exactly \(\rho\).
Reflection of the variable and value preserves the Taylor radius,
so \(g\) has this same radius.
\end{proof}

\subsection*{Evaluation of every conjectured choice}
In Theorem~\ref{counter}, equation \eqref{tau} is
\[
 1-\tau-\frac{\tau}{2(1-\tau)}=0.
\]
Its unique root in \((0,1)\) is \(\tau=1/2\).
Lemma~\ref{positive} therefore gives
\[
 R=F(1/2)=\frac12\sqrt{\frac12}\,e^{-1/2}.
\]
This is a radius for the entire complex Taylor disk, not just an
endpoint of real continuation.

On the other hand, the source saddle equation \eqref{saddles} becomes
\[
 (1-\phi)(1/2+\phi)-\phi=0
 \quad\Longleftrightarrow\quad
 (2\phi-1)(\phi+1)=0.
\]
Hence \(\phi=1/2\) or \(\phi=-1\); there are no further candidates.
Both are real, so \(|e^{2\pi i l\phi}|=1\) for every integer \(l\).
Substituting them into the remaining factors of \eqref{claim} gives
\begin{align*}
 \phi=\tfrac12:\quad
 R_+&=e^{-1}
 \left|(\tfrac12)^{1/2}(-1)^1
             (\tfrac12)^{1/2}(\tfrac12)^{1/2}\right|
       =\frac{e^{-1}}{2\sqrt2},\\
 \phi=-1:\quad
 R_-&=e^{-1}
 \left|2^2\,2^{-1/2}(\tfrac12)^{1/2}(-1)^{-1}\right|
       =\frac2e.
\end{align*}
The powers here are all defined, and changing any logarithm branch
does not affect their moduli because their exponents are real.
Finally \(R/R_+=e^{1/2}>1\), while
\[
 R<\frac12<\frac2e=R_-
\]
since \(1<e<3<4\). Neither saddle works for any integer phase.
This proves Theorem~\ref{counter}. \qed

\subsection*{The missing factor on the positive real class}
The broader lemma also clarifies the discrepancy. For its parameters,
put \(t_j=-a_j\) and
\[
 \phi_j=\frac{p_j\tau}{a_j-\tau}>0,\qquad
 \sum_j\phi_j=1-\tau.
\]
These numbers solve \eqref{saddles}. Taking absolute values in
\eqref{claim}, all integer phase factors disappear, and
\[
 \left|
 (1-\textstyle\sum_j\phi_j)^{1-\sum_j\phi_j}
 \prod_j
 \left(\frac{t_j}{p_j+\phi_j}\right)^{p_j+\phi_j}
 p_j^{p_j}\phi_j^{\phi_j}
 \right|
 =\tau\prod_j(a_j-\tau)^{p_j}.
\]
Indeed \(p_j+\phi_j=p_ja_j/(a_j-\tau)\) and
\((a_j-\tau)\phi_j/p_j=\tau\), so the powers of \(\tau\)
sum to \(\tau+\sum_j\phi_j=1\).
The correct positive-class value uses \(e^{-\tau}\), whereas the
printed conjecture uses \(e^{-1}\). Its ratio to the true value
is \(e^{\tau-1}\). Replacing the exponential by
\(e^{\sum_j\phi_j-1}\) repairs this particular real saddle.
No assertion about the repaired formula for arbitrary complex
parameters is made.

This counterexample satisfies the primary source's hypotheses,
uses a regular inverse germ with \(f'(0)=1\), and tests every
allowed saddle and integer phase. It is not based on a database
transcription or on a singular choice of input parameters.

\section*{Task 2: the corrected complex-parameter conjecture}
Insert the factor \(e^{\sum_j\phi_j}\) into the proposed expression.
Determine whether a critical-value selection rule and logarithm
branches give the inverse radius for arbitrary complex \(p_j,t_j\),
or construct a counterexample to that revised assertion.

\section*{Task 3: coefficient asymptotics in the positive class}
For Lemma~\ref{positive}, establish the dominant square-root expansion
at \(\rho\), determine whether any other point on \(|w|=\rho\)
is singular, and derive an explicit leading asymptotic for \(c_n\)
with a controlled remainder.

\section*{Task 4: sensitivity of the exact radius}
Differentiate the implicit equation for \(\tau\) and formula
\eqref{radius} with respect to each \(a_j,p_j\). Determine sharp
monotonicity and convexity properties, including the limits as
one exponent tends to zero or a negative root approaches zero.

\section*{Task 5: a certified critical-value algorithm}
For rational positive \(a_j,p_j\), construct an interval algorithm
that isolates the unique \(\tau\), encloses \(\rho\) to prescribed
accuracy, and gives a rigorous tail bound for the inverse Taylor
series throughout \(|w|\le r<\rho\).

\begin{thebibliography}{9}
\bibitem{Castle}
P. Castle, \emph{Taylor series for generalized Lambert W functions},
arXiv:1801.09904v1 (2018), 15 pp.; Conjecture 1.2 on p.~3
and its heuristic discussion in Section 6.
\url{https://arxiv.org/abs/1801.09904v1}.
\bibitem{Neto}
A. Francisco Neto,
\emph{New approach to the Lagrange--B\"urmann theorem via omega
calculus and applications},
Electron. J. Differential Equations \textbf{2025}, no.~90, 1--32.
\url{https://doi.org/10.58997/ejde.2025.90}.
Related literature: its Section 7 uses finite polynomial/rational
factors, and its Section 8 addresses the different questions of
Mez\H{o}--Baricz. It does not state a resolution of Castle's
arbitrary-power conjecture.
\end{thebibliography}
\end{document}
